\documentclass[../../../include/open-logic-section]{subfiles}

\begin{document}

\olfileid{his}{set}{infinitesimals}
\olsection{الكميات المتناهية الصغر والاشتقاق}

انتهى نيوتن ولايبنتس، كل واحد منهما على استقلال، إلى حساب التفاضل والتكامل
في أواخر القرن السابع عشر. ومن أجلّ ما استُعمل فيه هذا الحساب
\emph{الاشتقاق}؛ ومقصوده، على وجه التقريب، أن يُحدّ «معدل تغير» الدالة
عند نقطة، أي ميلها.

ولنجل الفكرة بمثال محسوس: تأمل الدالة
${\OLId{latin-ordinary}{f}}({\OLId{latin-ordinary}{x}}) = \nicefrac{{\OLId{latin-ordinary}{x}}^{\OLMathNumeral{2}}}{{\OLMathNumeral{4}}} + \nicefrac{{\OLMathNumeral{1}}}{{\OLMathNumeral{2}}}$، وقد رُسمت فيما يأتي
بالسواد:
\begin{center}
	\begin{tikzpicture}[scale=1]
	\draw[->, gray] (-1,0) -- (4.25,0) node[right] {${\OLId{latin-ordinary}{x}}$};
	\draw[->, gray] (0,-0.25) -- (0,5.25) node[above] {${\OLId{latin-ordinary}{f}}({\OLId{latin-ordinary}{x}})$};
	
	\foreach \x/\xtext in {1/1, 2/2, 3/3, 4/4}
	\draw[shift={(\x,0)}, gray] (0pt,2pt) -- (0pt,-2pt) node[below] {$\xtext$};
	
	\foreach \y/\ytext in {1/1, 2/2, 3/3, 4/4, 5/5}
	\draw[shift={(0,\y)}, gray] (2pt,0pt) -- (-2pt,0pt) node[left] {$\ytext$};
	
	\draw[oldiagcolorC] (0.5, 9/16) -- (3.5, 57/16) -- (3.5, 9/16)--cycle;
	\draw[oldiagcolorD] (.5, 9/16) -- (2.5, 33/16) -- (2.5, 9/16) -- cycle;
	\draw[oldiagcolorE] (.5, 9/16) -- (1.5, 17/16) -- (1.5, 9/16) -- cycle;
	\draw[thick] (-1,.75) parabola bend (0,0.5) (4,4.5);
	\end{tikzpicture}
\end{center}
ولنفرض أنا نريد ميل الدالة عند~${\OLId{latin-ordinary}{c}} = \nicefrac{{\OLMathNumeral{1}}}{{\OLMathNumeral{2}}}$. فنرسم أولاً مثلثاً
يقارب ميل وتره الميل عند تلك النقطة، وليكن هو المثلث الأحمر المتقدم. فإذا
كان~$\beta$ طول قاعدته، كان ارتفاعه
${\OLId{latin-ordinary}{f}}(\nicefrac{{\OLMathNumeral{1}}}{{\OLMathNumeral{2}}}+\beta) - {\OLId{latin-ordinary}{f}}(\nicefrac{{\OLMathNumeral{1}}}{{\OLMathNumeral{2}}})$؛ ومن ثم يكون ميل الوتر:
\[
\frac{{\OLId{latin-ordinary}{f}}(\nicefrac{{\OLMathNumeral{1}}}{{\OLMathNumeral{2}}}+\beta) - {\OLId{latin-ordinary}{f}}(\nicefrac{{\OLMathNumeral{1}}}{{\OLMathNumeral{2}}})}{\beta}.
\]
فمثلثنا~!!{colorC}، وطول قاعدته~${\OLMathNumeral{3}}$، يكون ميله مساوياً لـ${\OLMathNumeral{1}}$ سواء بسواء.
فإن أخذنا أصغر منه، أعني المثلث~!!{colorD} الذي طول قاعدته~${\OLMathNumeral{2}}$، كان وتره
أدنى إلى المطلوب، وميله~$\nicefrac{{\OLMathNumeral{3}}}{{\OLMathNumeral{4}}}$. ثم المثلث~!!{colorE} أصغر
منه أيضاً، وطول قاعدته~${\OLMathNumeral{1}}$، فيعطينا تقريباً أحسن، وميله
~$\nicefrac{{\OLMathNumeral{1}}}{{\OLMathNumeral{2}}}$.

وكلما صغّرنا المثلث ازداد تقريبنا جودة. فلنا أن نقول: إن وتر مثلث قاعدته
\emph{متناهية في الصغر} يعطينا الميل عينه عند
${\OLId{latin-ordinary}{c}} = \nicefrac{{\OLMathNumeral{1}}}{{\OLMathNumeral{2}}}$. وبهذا نحصل على صيغة المشتقة الأولى للدالة~${\OLId{latin-ordinary}{f}}$
عند النقطة~${\OLId{latin-ordinary}{c}}$:
\[
{{\OLId{latin-ordinary}{f}}'}({\OLId{latin-ordinary}{c}}) = \frac{{\OLId{latin-ordinary}{f}}({\OLId{latin-ordinary}{c}}+\beta) - {\OLId{latin-ordinary}{f}}({\OLId{latin-ordinary}{c}})}{\beta} \text{ حيث $\beta$ متناهية في الصغر.}
\]
وهذا، في الجملة، قول نيوتن ولايبنتس.

ولكن ما هذه الكمية \emph{المتناهية في الصغر}؟ ها هنا موضع الإشكال؛ فإن
كان~$\beta = {\OLMathNumeral{0}}$ لم تكن~${\OLId{latin-ordinary}{f}}'$ محدودة حداً صحيحاً، لاشتمالها على القسمة
على~${\OLMathNumeral{0}}$. وإن كان~$\beta > {\OLMathNumeral{0}}$ لم نصب الميل نفسه، وإنما أصبنا
\emph{تقريباً} منه.

وليس هذا اعتراضاً حادثاً أُلصق بالماضين؛ فهذا بيركلي يخاطب أتباع نيوتن
منتقداً:
\begin{quote}
I admit that signs may be made to denote either any thing or nothing:
and consequently that in the original notation ${\OLId{latin-ordinary}{c}} + \beta$, $\beta$
might have signified either an increment or nothing. But then which of
these soever you make it signify, you must argue consistently with
such its signification, and not proceed upon a double meaning: Which
to do were a manifest sophism. (\citealt[\S{}XIII]{Berkeley1734},
variables changed to match preceding text)
\end{quote}
وقد يدفع مدافع عن حساب المتناهيات في الصغر حجة بيركلي بأن القول بها مجاز
لا حقيقة؛ أي إنه متى أُخذ مثلث \emph{صغير جداً} حصل تقريب \emph{واف
بالغرض} للمماس. ولم يخل بيركلي من جواب عن هذا أيضاً: فلئن كفى في الهندسة،
فإنه ينقض \emph{مكانة} الرياضيات؛ إذ
\begin{quote}
we are told that \emph{in rebus mathematicis errores qu\`{a}m minimi
non sunt contemnendi}. [In the case of mathematics, the smallest
errors are not to be neglected.] \citep[\S{}IX]{Berkeley1734}
\end{quote}
والكلام المائل منقول، أو كاد، من كتاب نيوتن~\emph{تربيع المنحنيات}
(\OLClassicalDigits{1704}).

وكان لاعتراض بيركلي الفلسفي غور ونفاذ؛ على أن حساب التفاضل والتكامل نجح
نجاحاً عظيماً، ومضى الرياضيون يستعملونه من غير أن يوقعهم في خطأ.

\end{document}
